% APS and AIP journals with REVTeX 4.2. Change `prl` to prd, prb, pra, pre,
% prx, rmp or another journal, and `aps` to `aip` for AIP journals.
\documentclass[aps,prl,reprint,amsmath,amssymb]{revtex4-2}

\usepackage{graphicx}

\begin{document}

\title{An Example Title}

\author{A. U. Thor}
\email{author@example.org}
\affiliation{Department of Physics, Your University, City, Country}

\author{Second Author}
\affiliation{Another Institute, City, Country}

\date{\today}

\begin{abstract}
State the question, the method and the main result. Physical Review Letters
limits the abstract to 600 characters.
\end{abstract}

\maketitle

\section{Introduction}

Say what is known, what is missing, and what this paper adds. Cite earlier
work by key \cite{einstein1905,knuth1984}.

\section{Method}

Describe the setup and the analysis in enough detail to repeat them.

\begin{equation}
  E^2 = (pc)^2 + (mc^2)^2
  \label{eq:energy}
\end{equation}

\begin{table}[b]
  \caption{\label{tab:results}Example measurements.}
  \begin{ruledtabular}
    \begin{tabular}{lcc}
      Sample & Value & Uncertainty \\
      \hline
      A & 1.02 & 0.11 \\
      B & 0.97 & 0.08 \\
    \end{tabular}
  \end{ruledtabular}
\end{table}

\section{Results}

Report what you found. Table~\ref{tab:results} lists the measurements and
Eq.~(\ref{eq:energy}) is referenced like this.

\begin{acknowledgments}
Thank the people and grants that made the work possible.
\end{acknowledgments}

\bibliography{references}

\end{document}
